\itemhead{习题1.2 · 第2题（教材p.23）}{XJ-B1-C01-HW-P023-Y0993}
\label{XJ-B1-C01-HW-P023-Y0993}
{\footnotesize 湘教版教材p.23 / PDF第30页 / 习题1.2 / 第2题}
\field{原题}{写出下列命题的否定，并判断其真假。
\subq{1}$p$：边长分别为 5，12，13 的三角形是直角三角形；
\subq{2}$q$：面积相等的三角形都是全等三角形；
\subq{3}$s$：二次函数 $y=x^2-2x+3$ 的图象与 $x$ 轴没有交点。}
\field{参考答案}{（1）边长为5、12、13的三角形不是直角三角形，假；（2）存在面积相等但不全等的两个三角形，真；（3）$y=x^2-2x+3$ 的图象与 $x$ 轴有交点，假。答案册此处第三问的题干不同，已按本题重算。}
\field{核对说明}{版本题干变化：答案册PDF7第2题第3问为xy\ensuremath{\ne}0的否定；现教材问二次函数与x轴的交点。}
{\small \field{关联知识}{L01 命题、真假、逆命题与否定（教材帮p.30 / PDF39）；L04 含一个量词命题的否定（教材帮p.39、40 / PDF48、49）}
\field{方法与易错点}{明确论域与断言；真命题给理由，假命题给满足题设而使结论失败的反例。 易错：把逆命题当成否定；只改变结论中的个别词而改变原逻辑范围。}
\field{答案依据}{答案册相关题组：PDF第7页（册内页6）；按当前题库原题逐项重算，不按旧版题号机械对套。}}
